% =========================================================================
% CHAPTER 4: IMPLEMENTATION DETAILS
% =========================================================================
This chapter details the actual implementation of the proposed system. It includes information about the software tools, database tables, API routes, and critical code snippets.

\section{Technology Stack}
List and describe the programming languages, libraries, frameworks, and operating systems used. For example:
\begin{description}
    \item[Backend:] Python with FastAPI or ASP.NET Core for writing high-performance REST APIs.
    \item[Frontend:] TypeScript and React/Next.js for building a responsive, component-based user interface.
    \item[Database:] PostgreSQL or Firebase Firestore for storing structured application data.
    \item[Machine Learning:] PyTorch or TensorFlow for training and evaluating deep learning models.
\end{description}

\section{Algorithms and Pseudocode}
You can present the core operations of your system using pseudocode. Algorithm~\ref{alg:simple-process} shows an example using the \texttt{algorithm} package.

\begin{algorithm}[H]
\caption{Data Processing and Verification}
\label{alg:simple-process}
\begin{algorithmic}[1]
\REQUIRE Input dataset $D$, Threshold parameter $\tau$
\ENSURE Verified items list $V$
\STATE Initialize empty list $V$
\FORALL{item $x \in D$}
    \STATE $score \leftarrow \text{EvaluateModel}(x)$
    \IF{$score \ge \tau$}
        \STATE Add $x$ to $V$
        \STATE $\text{SendNotification}(x, \text{"Verified"})$
    \ELSE
        \STATE $\text{SendNotification}(x, \text{"Rejected"})$
    \ENDIF
\ENDFOR
\RETURN $V$
\end{algorithmic}
\end{algorithm}

\section{Code Listings}
You can include code listings directly in LaTeX. Listing~\ref{lst:python-sample} shows an example python function.

\begin{lstlisting}[language=Python, caption={Sample Python Inference Function}, label={lst:python-sample}]
import torch
import torch.nn as nn

def predict_sample(model, image_tensor, threshold=0.5):
    # Set model to evaluation mode
    model.eval()
    with torch.no_grad():
        output = model(image_tensor)
        probabilities = torch.sigmoid(output)
        prediction = (probabilities > threshold).float()
    return prediction.item(), probabilities.item()
\end{lstlisting}

Listing~\ref{lst:shell-sample} shows how to display a custom terminal/shell output using the \texttt{tcblisting} environment configured in the preamble.

\begin{tcblisting}{terminal shell, title={Server Execution Terminal}, label={lst:shell-sample}}
$ python main.py --mode=train --epochs=10
Loading training datasets...
Training started. Epoch 1/10 - loss: 0.652 - acc: 0.723
Training completed. Model saved to outputs/best_model.pth
\end{tcblisting}
